% GERADO POR lab/sumario.py — NÃO EDITE À MÃO.
% Os enunciados numerados do livro, na ordem em que aparecem: número, título e página
% saem do livro.aux da compilação anterior; a espécie sai do fonte do capítulo.

\noindent\hyperref[prop:barra]{\textbf{Proposição 1.1 --- a barra da fração}}\dotfill\ \pageref{prop:barra}\par
\noindent\hyperref[prop:passo]{\textbf{Proposição 2.1 --- o passo e o horizonte}}\dotfill\ \pageref{prop:passo}\par
\noindent\hyperref[prop:conta]{\textbf{Proposição 3.1 --- a conta do corte}}\dotfill\ \pageref{prop:conta}\par
\noindent\hyperref[prop:teto]{\textbf{Proposição 4.1 --- o teto independente}}\dotfill\ \pageref{prop:teto}\par
\noindent\hyperref[prop:piso]{\textbf{Proposição 4.2 --- o piso do atraso}}\dotfill\ \pageref{prop:piso}\par
\noindent\hyperref[prop:particao]{\textbf{Proposição 6.1 --- o preço da partição}}\dotfill\ \pageref{prop:particao}\par
\noindent\hyperref[prop:teto-do-par]{\textbf{Proposição 7.1 --- o teto do par}}\dotfill\ \pageref{prop:teto-do-par}\par
\noindent\hyperref[def:resposta]{\textbf{Definição 12.1 --- contenção e resposta}}\dotfill\ \pageref{def:resposta}\par
\noindent\hyperref[prop:piso-contencao]{\textbf{Proposição 12.2 --- o piso da contenção}}\dotfill\ \pageref{prop:piso-contencao}\par
\noindent\hyperref[def:escala-forma]{\textbf{Definição 13.1 --- escala e forma}}\dotfill\ \pageref{def:escala-forma}\par
\noindent\hyperref[prop:padronizacao]{\textbf{Proposição 13.2 --- a escala some na padronização}}\dotfill\ \pageref{prop:padronizacao}\par
\noindent\hyperref[def:esquecimento]{\textbf{Definição 14.1 --- esquecimento com taxa}}\dotfill\ \pageref{def:esquecimento}\par
\noindent\hyperref[prop:meia-vida]{\textbf{Proposição 14.2 --- a meia-vida e a janela}}\dotfill\ \pageref{prop:meia-vida}\par
\noindent\hyperref[prop:nulo-da-direcao]{\textbf{Proposição 15.1 --- o desvio do nulo}}\dotfill\ \pageref{prop:nulo-da-direcao}\par
\noindent\hyperref[def:episodio]{\textbf{Definição 16.1 --- episódio dirigido}}\dotfill\ \pageref{def:episodio}\par
\noindent\hyperref[prop:recorde]{\textbf{Proposição 17.1 --- o recorde e o que vem depois}}\dotfill\ \pageref{prop:recorde}\par
\noindent\hyperref[prop:reconstrucao]{\textbf{Proposição 18.1 --- o erro de cada reconstrução}}\dotfill\ \pageref{prop:reconstrucao}\par
\noindent\hyperref[prop:atraso-parado]{\textbf{Proposição 21.1 --- o atraso não cobra dias}}\dotfill\ \pageref{prop:atraso-parado}\par
